% | docs/security.tex
%
% The paper proof of the Midfall Halo2 verifier's knowledge soundness, and of the
% deployed Solidity contract's. Plan Section 10 step 1 says write this BEFORE the
% Lean proofs; it is also where the findings recorded in docs/REVIEW.md have to be
% resolved, because several of them are statements about the protocol that the
% Rust source cannot answer.
%
% STATUS: skeleton with the settled parts written out and every open question
% named. Environments marked OPEN are not claims.

\documentclass[11pt]{article}
\usepackage{amsmath,amssymb,amsthm}
\usepackage[margin=1in]{geometry}
\newtheorem{theorem}{Theorem}[section]
\newtheorem{lemma}[theorem]{Lemma}
\newtheorem{definition}[theorem]{Definition}
\newtheorem{assumption}[theorem]{Assumption}
\newtheorem{open}[theorem]{Open problem}
\newcommand{\Fr}{\mathbb{F}_r}
\newcommand{\Gone}{\mathbb{G}_1}
\newcommand{\Gtwo}{\mathbb{G}_2}
\newcommand{\GT}{\mathbb{G}_T}
\newcommand{\Hash}{\mathsf{H}}
\title{Knowledge soundness of the Midfall Halo2 verifier and of its deployed
Solidity lowering \\ \large (working draft)}
\author{Midfall formal project}
\date{Baseline \texttt{formal-baseline-v1}, Lean port at \texttt{formal/}}

\begin{document}
\maketitle

\begin{abstract}
We state and prove, in the algebraic group model with a random oracle, knowledge
soundness for the verifier implemented in \texttt{proofs/src} and ported to Lean in
\texttt{formal/Midfall/Halo2/Exec}, and for the Solidity/Yul contract generated from
a verifying key by \texttt{proofs/solidity-verifier}. The security argument is a
reduction to $q$-DLOG. The lowering argument is a semantic equivalence at the call
boundary between the Lean verifier and the compiled contract, for \emph{every} hash
function, which is what lets the same adversary bound transfer to the deployed
bytecode. Five things are assumed rather than proved: bilinearity and non-degeneracy
of the pairing as the ecosystem implements it, the EIP-2537 and EIP-198 precompile
semantics, $q$-DLOG, that the imported verifying-key representative binds the key,
and that the Lean port is faithful to the Rust. Section~\ref{sec:open} lists what is
not yet established, including two findings that currently make statements in the
project plan unprovable as written.
\end{abstract}

\section{The protocol}

One round, single-phase. A constraint system $\mathcal{C}$ over $\Fr$ with $n=2^k$
rows, advice columns, instance columns, fixed columns, a permutation argument,
\texttt{logUp} arguments and trash arguments. The prover commits, the verifier
squeezes challenges from a Keccak hash chain, and the final test is one pairing
equation.

\begin{definition}[transcript]\label{def:transcript}
The transcript is a sponge over byte strings: absorb concatenates into a buffer,
squeeze is $\Hash(\mathit{state})$, after which the state \emph{becomes} that
digest (\texttt{transcript/mod.rs:139-141}). A challenge is the \emph{big-endian}
digest reduced mod $r$ (\texttt{implementors.rs:20-39}). The wire encoding of a
scalar is little-endian while its transcript encoding is big-endian, and a
non-canonical repr is rejected (rule R7). Both are pinned by known-answer tests in
the port, after finding R-20 showed that a self-consistent test cannot detect a
byte-order error.
\end{definition}

\begin{definition}[rounds]\label{def:rounds}
In order: $\theta$ after the advice commitments; $\beta,\gamma$ after the lookup
multiplicities; $\mathit{trash}$ after the permutation and lookup commitments; $y$
after the trash commitments; then the quotient limb commitments; then $x$. The
multi-open challenges $x_1,x_2,x_3,x_4$ are squeezed inside
\texttt{multi\_prepare}. This order is the ported \texttt{parse\_trace} plus
\texttt{verify\_algebraic\_constraints}, and it is \emph{protocol}: reordering two
squeezes changes the language the verifier accepts, not merely its output.
\end{definition}

\section{Assumptions}

\begin{assumption}[$q$-DLOG]\label{as:qdlog}
For the SRS size $N$ in use, $q$-DLOG is hard in $\Gone$: given
$\{\tau^i G_1\}_{i<N}$, no ppt adversary outputs $\tau$ with non-negligible
probability. ArkLib's \texttt{Groups.arsdhAssumption} is the closest library
statement; whether we reduce to it or state our own is a Phase-4 choice.
\end{assumption}

\begin{assumption}[pairing]\label{as:pairing}
$e : \Gone \times \Gtwo \to \GT$ is bilinear and non-degenerate. We do not prove
this of the ecosystem implementation: \texttt{Assumptions/Pairing.lean} carries it
as a named class \texttt{AdmissiblePairing} with the hypotheses as fields, and the
concrete BLS12-381 instance is validated differentially rather than proved (D-P1).
\end{assumption}

\begin{assumption}[key binding]\label{as:vk}
The verifying key's transcript representative $\mathit{repr}$ is the image of the
key under a collision-resistant function of the constraint system and commitments.
The port \emph{imports} this scalar: recomputing it requires byte-exact Rust
debug formatting for two types, which the upstream code itself flags as provisional
(finding R-1). Everything downstream is bound to it, so the binding is an assumption
about the importer, checked by conformance rather than proved.
\end{assumption}

\begin{assumption}[precompiles]\label{as:precompiles}
EIP-2537 \texttt{G1ADD}, \texttt{G1MSM}, \texttt{PAIRING\_CHECK} and EIP-198/7883
\texttt{MODEXP} behave as the EIP texts specify, byte for byte.
\texttt{Yul/Externals.lean} states the relations; they are validated against revm in
Phase 5. Two EIP subtleties the port records: \texttt{G1ADD} performs no subgroup
check, and \texttt{G1MSM} does not reduce scalars mod $r$.
\end{assumption}

\section{The per-round bad-point analysis}\label{sec:rounds}

Each challenge $c$ is squeezed after the commitments that could depend on it, so
$c$ is uniform in $\Fr$ given the transcript. For a batched identity
$P(c) = \sum_i c^i P_i$, the standard step is: if the $P_i$ are not all zero as
polynomials then $P$ has at most $\deg P$ roots, so
$\Pr_c[P(c)=0] \le \deg P / r \le 3\deg P/2^{256}$. The factor $3$ is
$\lceil 2^{256}/r\rceil$, which is \emph{proved} in \texttt{Spec/Degrees.lean}
from CompPoly's primality certificate for $r$.

\begin{open}[nonzeroness]\label{op:nonzero}
For each round we must show the batched polynomial is not identically zero, given
that the prover's commitments are not representations of a satisfying witness. This
is the actual content of the proof and it is not a fact about the Rust code. The
counts in \texttt{Spec/Degrees.lean} charge each identity its full degree and sum,
which is sound \emph{if} this holds. Sources: PLONK for the gate and permutation
parts, the logUp paper for the lookup part, and the trash argument, which is ours.
\end{open}

\subsection{Where the standard argument currently fails}\label{sec:truncation}

Finding R-22. The deployed profile enables \texttt{truncated-challenges}, and the
truncation is applied \emph{per round}, not globally: $\theta,\beta,\gamma$,
$\mathit{trash}$, $y$ and $x$ are full $\Hash \bmod r$, while the multi-open
batching coefficients are $\mathsf{trunc}(x_1^j)$ and $\mathsf{trunc}(x_4^j)$
where $\mathsf{trunc}$ keeps the low 128 bits of the canonical representative
(\texttt{arithmetic.rs:247-257}, \texttt{kzg/mod.rs:393,501,520}). The batching
equation is then
\[ \sum_j \mathsf{trunc}(x_1^j)\, C_j = 0 \quad\text{in } \Fr, \]
which is \emph{not} a polynomial equation in $x_1$. No degree bound applies, so
$\mathsf{roundErr}(x_1)$ and $\mathsf{roundErr}(x_4)$ as written in the plan are
unproven for the configuration we deploy.

\begin{open}[truncated batching]\label{op:trunc}
Bound, for a fixed nonzero tuple $(C_j)$ with known representations, the number of
$x_1 \in \Fr$ with $\sum_j \mathsf{trunc}(x_1^j) C_j = 0$. Heuristically the
truncation leaves about $2^{128}$ bits of effective randomness per test, which
would give a bound of shape $m/2^{128}$ for a different $m$ than the maximum number
of polynomials in a point set; proving it needs to relate the integer map
$x \mapsto x^j \bmod r \bmod 2^{128}$ to polynomial structure. Options: prove
this; drop \texttt{truncated-challenges} from the supported profile; or state the
theorems only for \texttt{truncatedChallenges = false}, which would exclude the
deployment.
\end{open}

Two smaller corrections the same analysis forces (R-23, now executed by D-P21 in
\texttt{Security/Challenges.lean} and \texttt{Spec/Degrees.lean}). Only $x_3$ has
its \emph{value} truncated, and the order is reduce-then-truncate
(\texttt{kzg/mod.rs:455-457}: \texttt{squeeze\_challenge} and then
\texttt{truncate}), so the two factors \emph{multiply} rather than replace. A
truncated value has $\lceil r/2^{128}\rceil$ field preimages, each hit at most
$\lceil 2^{256}/r\rceil = 3$ times out of $2^{256}$, so the per-value mass is
$3\lceil r/2^{128}\rceil/2^{256} = 1.3585351690\ldots/2^{128}$, and $s/2^{128}$
is \emph{not} an upper bound for $x_3$. The earlier estimate of $1.5$ overstated
the shortfall: it is $0.442$ bits of budget, not half a bit. For the six rounds
whose value is not truncated the uniform $s/2^{128}$ was sound but $2^{126.42}$
too large, which is why $\mathsf{chalMass}$ now takes the round. Separately,
$x_3$ must avoid both the $n$-th roots of unity and every distinct opening point,
so the honest count is $n + \sum_i |S_i|$; the plan's $\max_i |S_i| \cdot n$
neither dominates it nor follows from it, so the row charges
$\max(\max_i |S_i| \cdot n,\; n + \sum_i |S_i|)$ until the paper decides which
half is real.

\section{The verifier's final test}

Without an accumulator the verifier accepts iff
$e(L, sG_2)\cdot e(R, -G_2) = 1$ for the DualMSM $(R,L)$ built by
\texttt{multi\_prepare}. With one, the contract batches: accept iff
$e(R+\alpha R', G_2)\cdot e(L+\alpha L', -sG_2) = 1$ where
$\alpha = \Hash(\mathit{tag} \,\|\, \mathit{repr} \,\|\, R \,\|\, L
\,\|\, R' \,\|\, L') \bmod r$ with $0 \mapsto 1$.

\begin{lemma}[batching completeness]\label{le:batchcomp}
If both pairing products are $1$, the batched equation holds for every $\alpha$.
\end{lemma}
\begin{proof}
Bilinearity (Assumption~\ref{as:pairing}) gives
$e(R+\alpha R',G_2)\,e(L+\alpha L',-sG_2) =
[e(R,G_2)\,e(L,-sG_2)]\cdot[e(R',G_2)\,e(L',-sG_2)]^{\alpha} = 1$.
\end{proof}

\begin{lemma}[at most one bad $\alpha$]\label{le:batchroot}
If the two separate checks do not both pass, the set of $\alpha$ for which the
batched check passes has at most one element.
\end{lemma}
\begin{proof}
Write $X = e(R,G_2)\,e(L,-sG_2)$ and $Y = e(R',G_2)\,e(L',-sG_2)$. The batched
test is $X Y^{\alpha} = 1$. If $Y = 1$ then $X \ne 1$, else both checks pass, and
there is no solution; if $Y \ne 1$ there is exactly one solution.
\end{proof}

\begin{theorem}[$\alpha$-batching soundness]\label{th:alpha}
For an adversary making $q_{\Hash}$ oracle queries, the probability that the
batched check passes while the two separate checks do not both pass is at most
$6(q_{\Hash}+1)/2^{256}$.
\end{theorem}
\begin{proof}[Proof sketch]
By Lemma~\ref{le:batchroot} success requires guessing the unique $\alpha$, and
$\alpha$ is the image of a transcript containing the pair the adversary must commit
to before $\alpha$ exists; each oracle query wins with probability at most
$\lceil 2^{256}/r\rceil$ times the number of group elements whose relation to
$\alpha$ must be consistent. Formalising the commit-before step is the lazy-sampling
argument in \texttt{Security/ROM}.
\end{proof}

\begin{open}[Rust versus contract in accumulator mode]\label{op:acc}
Findings R-10 and R-26: the Rust \texttt{verify\_final} batches with a scalar
derived from a POW2 sponge over a different encoding, the contract uses Keccak with
the key representative, the two order the four points differently, and the
contract's decoder assumes a \emph{collapsed} accumulator shape that the Rust
encoder does not guarantee. So the claim that the contract checks the same thing as
Rust is true only relative to the contract's own reference predicate
\texttt{ContractRef.accepts}, which is exactly how plan Section 9.5 states it. The
exact-equivalence theorem is therefore the no-accumulator one.
\end{open}

\section{Main theorem}\label{sec:main}

\begin{theorem}[S1, conditional]\label{th:s1}
Assume \ref{as:qdlog}--\ref{as:precompiles} and that Open
problems~\ref{op:nonzero} and~\ref{op:trunc} resolve with the bounds stated in
\texttt{Spec/Degrees.lean}. Then for a single-phase well-formed key the verifier's
knowledge error is at most $\mathsf{ksErr}(\mathcal{C}, q_{\Hash}) +
\Pr[\mathsf{qdlog}\text{ game wins}]$, with a straight-line extractor reading the
advice witness off the AGM representations of the advice commitments.
\end{theorem}

\section{The lowering}\label{sec:lowering}

Section 9.5 of the plan states the headline: for every supported key and config,
every calldata, and \emph{any} hash $\Hash$, the compiled contract returns ABI
\texttt{true} iff the Lean verifier accepts the decoded calldata, and otherwise
reverts or halts without changing state. Three properties make that the right
statement.
\begin{itemize}
\item It is \emph{observational}: we do not claim memory equality, only the outcome
at the call boundary. Lean error classes all map to revert.
\item It is \emph{hash-agnostic}, which is what lets the security theorem reuse it
with $\Hash$ instantiated as the random oracle.
\item It concerns the \emph{compiled} contract, i.e.\ solc's IR of the generated
Solidity file, translation-validated by \texttt{Yul/IRCheck.lean}, not merely the
assembly blocks we wrote. The Yul optimiser and the bytecode step stay trusted and
are covered by revm differential tests.
\end{itemize}

\begin{open}[public-input malleability]\label{op:mal}
Finding R-25 and the \texttt{calldata\_canonical} risk: the contract's accumulator
decoder ignores the coordinate words of an identity-flagged point
(\texttt{AccumulatorHelpers.yul:130-141}), so distinct instance tails decode to the
same accumulator. This is malleability, not forgery, since the decoded point is the
identity either way, but it falsifies injectivity as stated, and the fix is to
normalise in \texttt{accWords} rather than in the decoder.
\end{open}

\section{What is assumed about the port}\label{sec:port}

The Lean L0 is \emph{validated} against the Rust, not proved equivalent: Rust is not
a formal object. The validation is (i) line-faithful porting with provenance
comments, (ii) a second-agent review of every module, (iii) a conformance harness
replaying every absorb, squeeze, evaluation and MSM term, and (iv) a mutation
corpus. Where the Rust can panic, the port returns a named fault at the same site,
so a malformed key fails identically rather than diverging. Findings R-2 (three
faces), R-7, R-13, R-14, R-16, R-19, R-20 and R-25 are all products of that process,
and R-20 in particular shows that a test which compares a function against itself
proves nothing about orientation.

\section{Open problems, collected}\label{sec:open}
\begin{enumerate}
\item \ref{op:nonzero}: nonzeroness of each batched identity, the real content.
\item \ref{op:trunc}: the truncated batching coefficients, currently a gap in S1/S3
for the deployed profile (R-22).
\item \ref{op:acc}: the Rust/contract accumulator divergence (R-10, R-26).
\item \ref{op:mal}: public-input malleability and \texttt{calldata\_canonical}
(R-25).
\item The $\delta$ hypothesis: nothing checks that $\mathsf{DELTA}$ has order at
least the number of permutation columns, so the two permutation products can
coincide identically and the $\beta,\gamma$ degree argument collapses (R-3
extended).
\item \texttt{srsBound}: the plan's $2^{\mathit{extended\_k}}$ may be too small if
the proof needs $n(c+1)$ (R-23).
\item The Yul optimiser and bytecode hop stays trusted (plan Section 13).
\end{enumerate}

% Each entry below was checked against the primary source page, not recalled. Two
% entries were wrong before: plookup's title was invented, and the logUp attribution
% named a paper that is not findable under that title. See docs/plans/PROOF-PLAN.md,
% section "Bibliography corrections", for what each was corrected to and why.
\begin{thebibliography}{9}
\bibitem{plonk} Gabizon, Williamson, Ciobotaru. PLONK: Permutations over
Lagrange-bases for Oecumenical Noninteractive arguments of Knowledge.
Cryptology ePrint Archive 2019/953, 2019.
\bibitem{plookup} Gabizon, Williamson. plookup: A simplified polynomial protocol
for lookup tables. Cryptology ePrint Archive 2020/315, 2020.
\bibitem{logup} Hab\"ock. Multivariate lookups based on logarithmic derivatives.
Cryptology ePrint Archive 2022/1530, 2022. This is the formal reference for the
logarithmic-derivative (logUp) argument used by this codebase.
\bibitem{logupchar} Eagen and Hab\"ock. Bypassing the characteristic bound in
logUp. Cryptology ePrint Archive 2024/2067, 2024. Directly relevant to the
$\operatorname{char} > n$ hypothesis carried in Section~\ref{sec:rounds}: they
remove that bound by abstracting pole multiplicity. Whether their construction
applies to the trash/logUp variant implemented here is not yet established.
\bibitem{logupstar} Soukhanov. Logup*: faster, cheaper logup argument for
small-table indexed lookups. Cryptology ePrint Archive 2025/946, 2025.
\bibitem{halobook} The Halo 2 book: permutation argument, quotient splitting,
multi-open.
\bibitem{agm} Katz, Krenn, Thaler, Zaveri. On the Source of Statistical Parameters,
or How to Rework Old Proofs in the Ideal Model with Minimal Effort.
\bibitem{eip2537} EIP-2537: Precompile for BLS12-381 curve operations.
\bibitem{eip7883} EIP-7883: Modexp gas cost repricing.
\end{thebibliography}

\end{document}
